\clearpage
\section*{Concluzii}
\addcontentsline{toc}{section}{Concluzii}

Jocul elaborat a fost finisat cu succes deși mai necesită unele completăriă. În
timpul dezvoltării jocului am studiat despre modul cum sunt create jocurile
și am văzut care sunt dificultățile care pot apărea pe parcursul dezvoltării unui joc. Am început 
să realizez cât de dificilă poate fi dezvoltarea unui joc de proporțiii mari și
am descoperit o pagină web \cite{GameIdea} care aproximează valoarea unui joc
cu mecanici asemănătoare jocului meu la 7000 de \$. Realizarea acestui proiect
mi-a schimbat viziunea despre nivelul de dificultate al creării unui joc. În
cadrul acestei practici am avut ocazia să utilizez limbajul Rust care este
limbajul meu preferat și să îmi îmbunătățesc cunoștințele despre acesta.

Pentru a realiza acest proiect a trebuit să învăț Bevy și Rapier cu, care nu am
mai avut de aface înainte de acest moment. Utilizarea lor a fost destul de
dificilă și consider că dacă aș fi utilizat Unity aș fi făcut un joc mai
complet fiind că există un număr mai mare de tutoriale video pe YouTube în
comparație cu Bevy care nu are nici un tutorial pentru jocuri 3D. Totodată o
mare problemă este faptul că Bevy nu dispune de un editor de scenă ceea ce face
dezvoltarea jocurilor mult mai grea în comparație cu Unity.

Scopul jocului dezvoltat consă în colectarea de portocale care dau un scor.
Timpul este limitat, jucătorul având la dispoziție 60 de secunde pentru a
acumula un scor cât mai mare. Pentru viitor îmi propun să generez procedural 
plantele și să implementez un meniu. Totodată ar fi bine să optimizez codul și
să distribui acest joc pe \href{https://itch.io/}{Itch.io} ca să pot afla
părerea altor persoane despre el.

Datorită acestei experiențe am o înțelegere de bază a modului de creare a jocurilor
și a arhitecturii ECS astfel trecerea la Unity și C\# nu va fi dificilă.
Numărul mare de tutoriale face atractivă trecerea la Unity cel puțin pentru o perioadă pentru a însuși mai rapid modul de dezvoltarea al jocurilor. Totuși aș prefera să creez jocuri utilizând Rust și sper că în versiunile următoare Bevy va oferi un editor de scenă.

% Motivul pentru care voi utiliza Unity pentru o perioadă este numărul mare de tutoriale care ar permite o însușire rapidă a conceptelor din dezvoltarea jocurilor video, apoi adată ce voi avea o înțelegere mai bună a procesului de dezvoltare de jocuri și odată ce Bevy va dispune de un editor al scenei mă voi reîntoarce la Rust.
% Elaborând aceast proiect am studiat modul de dezvoltare a jocurilor video și am dezvoltat un joc 3D în care jucătorul are ca obiectiv să strângă cât mai multe portocale în decursul a 60 de secunde. Am acumulat experiența necesară pentru a realiza niște jocuri simple 3D utilizând limajul Rust și motorul Bevy. Am realizat cât de dificil este de realizat un joc mare într-o echipă mică acest lucru fiind confirmat și de această pagină web\cite{GameIdea} care estimează că costul unui joc cu lumea deschisă este mai mare de 7000 \$.
% Deși jocul creat este finisat, adică are un obiectiv și este finit nu am realizat tot ce mi-am pus ca scop, de aceea există posibilitatea dezvoltării acestui proiect în viitor. Generarea procedurală a plantelor este aspectul la care aș fi vrut să lucrez cel mai mult însă din cauza faptului că nu am avut experiență în crearea jocurilor celelalte sarcini au durat prea mult timp. La moment jocul nu are o interfață de utilizator cee ce nu îl face foarte intuitiv, jucătorul fiind aruncat în lumea jocului fără nici o explicație, pentru următoare versie se merită de implementat această interfață de utilizator. În ultimul rând jocul nu a fost distribuit pe nici una dintre platformele enumerate în capitolul \ref{section:platforme} de aceea după cizelarea jocului acesta va fi distribuit pe platforma \href{https://itch.io/}{Itch.io}.
% Pentru următoarele jocuri pe care le voi dezvolta voi trece la Unity și C\# fiind că există multe tutoriale(în bevy nu există foarte multe, maie ales despre jocuri 3D) ceea ce ar permite o însușire mai rapidă a disciplinei de dezvoltare a jocurilor video. Și voi reveni la Rust odată ce voi avea mai multe experiență în dezvoltarea jocurilor video sau pentru utilizarea acestui limbaj de programare în alte domenii.
